% Architecture and workflow of ETHOS.TSAM — the manuscript's Fig. 1.
%
% ── V2: WHAT CHANGES AGAINST architecture_overview.tex ────────────────────
% One change only, so the two can be compared side by side: Algorithms grows
% from a name-only box into a component with a members compartment, and its
% bottom edge drops from -123 to -135 to align with Configuration. Everything
% else — the round-cornered phases, the activity name, the legend, the page
% size — is identical.
%
% WHAT IS LISTED, AND WHY THOSE. The rule is "the operations the Pipeline
% actually calls into", checked against the imports rather than against a
% reading of the package:
%     pipeline/clustering.py:13-15  assign_clusters, cluster_and_represent,
%                                   representations, deterministic_argmin
%     pipeline/segmentation.py:7    segmentation
% Of those, `cluster_and_represent` is the convenience wrapper around the other
% two and `deterministic_argmin` is a tie-break helper, so the compartment
% shows the three that the phase boxes NAME as stages — assign clusters,
% representations, segment — which is what makes the box readable as "this is
% where Phase 2 and Phase 3 do their work". The two exact-optimisation classes
% follow on an indented run-in line, in the same shape as Configuration's
% ".representation takes" line, because they are what the external edges are
% about: k_medoids_exact.py:14 is the ONLY module in the package that imports
% Pyomo, which is what the "exact kmedoids" edge label says.
%
% NOT listed, deliberately: the six cluster methods and six representations.
% aggregation_methods.tex already enumerates every one of them as chips
% (its lines 176-192), and repeating a companion figure's content here would
% be twelve names that say nothing new. `compute_ordering` is also absent —
% it is real, but concurrency.py is reached from the distribution
% representation, NOT from the pipeline, so listing it under this box's rule
% would have been false.
%
% VISIBILITY MARKERS. `+` is UML public and is what every other compartment in
% this figure uses. The Algorithms members take `~`, UML PACKAGE visibility:
% visible inside tsam, not exported through tsam.__all__. That is the accurate
% marker and it agrees with the box's white fill, which the legend already
% reads as "internal to tsam". `-` (private) would have been wrong — Pipeline
% calls these. It is set as $\sim$ and not \textasciitilde: under T1 the text
% tilde sits at diacritic height and reads as an accent, the math one sits on
% the axis and reads as a marker.
%
% ── NOTATION: UML 2.5.1 ───────────────────────────────────────────────────
%   * Component and its nested internal structure: clause 11, Structured
%     Classifiers. Pipeline is drawn as a WHITEBOX. That is also arc42 §5
%     Level 2 (the whitebox/blackbox recursion), not a deviation from it.
%   * Dependency (dashed line, open arrowhead): clause 7, Common Structure.
%   * Action (ROUND-CORNERED rectangle): clause 16, Actions.
%   * ObjectFlow (solid line, open arrowhead): clause 15, Activities.
%   * Activity name in the upper left of the region that frames it: clause 15.
%
% ── WHY THE PHASES ARE ROUND-CORNERED AND THE COMPONENTS ARE NOT ──────────
% The four phases are not components. They are the temporal steps of one run,
% and an earlier revision drew them as sharp-cornered boxes with the same
% border as Tools/API/Results, which said the opposite: that the pipeline is
% built out of four more parts of the same kind as the boxes above it.
%
% UML already separates structure from behaviour by SHAPE, so no invention is
% needed. A Classifier is a sharp-cornered rectangle; an Action is a
% round-cornered one (clause 16). Making the phases round-cornered is therefore
% the whole fix, and it costs nothing else: colour keeps its existing job
% (public API / internal / external for the components, amber for optional
% stages) and shape alone carries structure-versus-behaviour. It also survives
% a greyscale print, which a colour difference would not.
%
% Two consequences follow, and both make the figure MORE conformant than the
% revision before it:
%
%   * The chain between the phases is now an ObjectFlow — an activity edge,
%     solid with an open arrowhead (clause 15) — and its labels name the object
%     that flows. The previous revision had to draw it as a Connector between
%     nested parts (clause 11.4) and lean on §11.7 to let that Connector
%     realize an InformationFlow, purely to win a solid line that a dashed
%     Dependency could not be told apart from. That liberty is gone: an
%     ObjectFlow between Actions is solid because UML says it is.
%   * The Pipeline's second compartment is the behaviour it owns, so it carries
%     the activity's name in its upper left, \guillemotleft activity\guillemotright{} run\_pipeline
%     (src/tsam/pipeline/orchestrator.py:697). That is where UML puts an
%     Activity's name, and it tells the reader in one line that what follows is
%     a run, not a parts list.
%
% The legend gains one row for the round-cornered swatch. It is paid for by
% giving the middle legend row three columns instead of two, so the figure's
% height does not move — see the PRINT TARGET note, there was ~1.6 mm of slack
% and a fourth row needs 8.
%
% ── WHY CONFIGURATION AND COMMONS HAVE NO ARROWS ──────────────────────────
% Both are used by EVERY building block. An earlier revision drew one arrow
% into each (Pipeline -> Configuration, Algorithms -> Commons), which is worse
% than drawing none: it asserts a specific dependency that is not the real
% relationship, and a reader is entitled to conclude that nothing else uses
% them. UML has no "used by everything" edge, so the conformant device is a
% KEYWORD on the component: both carry \guillemotleft shared\guillemotright, both arrows are gone, and
% the legend says what the keyword means.
%
% Fill colour keeps its existing job (blue = exposes public API, white =
% internal, grey = external), so the two channels do not collide: Configuration
% stays blue because ClusterConfig and friends are public, Commons stays white.
%
% Each keyword is declared ONCE. \guillemotleft use\guillemotright and \guillemotleft shared\guillemotright live in the legend;
% \guillemotleft activity\guillemotright appears once in the drawing, on the compartment it names,
% because that is the position that carries its meaning.
%
% ── OPTIONAL STAGES ───────────────────────────────────────────────────────
% Stages that run only when the caller configures them are set in AMBER inside
% the phase boxes, which is the same optional-means-amber convention the other
% two figures of this manuscript already use. Verified against
% src/tsam/pipeline/orchestrator.py:
%   Phase 1  prepare_data          normalize, unstack_to_periods,
%                                  [weights], [period sums]
%   Phase 2  cluster_candidates    assign clusters, representations
%   Phase 3  refine_representatives [extremes], [unweight], count occurrences,
%                                  [rescale], [segment]
%   Phase 4  build_result          denormalize, reconstruct + accuracy
% Note `unweight` is conditional — orchestrator.py:514 guards it on
% `prepared.weight_vector is not None` — and `count occurrences` (line 519) is
% the only unconditional step in Phase 3. An earlier draft of this figure
% called the whole of Phase 3 optional, which was wrong.
%
% ── PRINT TARGET ──────────────────────────────────────────────────────────
% US Letter, 1 inch margins, text column 165.1 x 228.6 mm, every figure placed
% at full text width (wp:extent cx="5943600" EMU = 6.5 in).
%   * ASPECT h/w <= 1.24, or the scaled figure outgrows the page.
%   * WIDTH <= ~165 mm, or Word scales the drawing DOWN and every font with it.
% The drawing measures 162.1 x 199.4 mm, aspect 1.230 — about 1.6 mm of height
% to spare. Anything that adds a row has to take one away somewhere else.
%
% ── LAYOUT GRID (millimetres, y negative downwards) ───────────────────────
%   columns  c1 = 24 (2..46)   c2 = 77 (55..99)   c3 = 133 (109..157)
%   row A    y -16..-48    Tools (c1)  API (c2)  Results (c3)
%   Pipeline y -56..-97    whitebox, x 2..157; header -60, compartment rule -64
%     the activity name sits in the compartment's upper left at (11,-64.5),
%     anchored north west; it is \tiny, ~2.4 mm tall, so it clears the phase
%     tops at -68 by ~1.1 mm. That 4 mm band is the only space it has.
%     four parts 22 mm wide on a 41.67 mm pitch, 19.67 mm gaps, y -68..-93:
%       6..28 (c 17)          47.7..69.7 (c 58.7)
%       89.3..111.3 (c 100.3) 131..153 (c 142)
%     object flows run at y -80.5 across the gaps; their labels are
%     free-standing nodes anchored south at y -76.5, each inside its own
%     19.67 mm gap. A label node is 16 mm of text width plus 0.6 mm of inner
%     sep a side = 17.2 mm, so it clears the phase boxes by ~1.2 mm a side. The
%     first attempt used 24.45 mm parts and 16.4 mm gaps, where 17.2 mm did not
%     fit at all and "Refined Representatives" overlapped the Phase 4 box.
%     They are NOT `pos=` labels: pos on a multi-segment path applies to the
%     last segment only, and a white-filled label centred on the flow would
%     leave only a stub of visible line at each end.
%   row C    y -105..-135  Commons (c1, -105..-123, name only — it genuinely
%     has nothing to list)  Algorithms (c2, -105..-135)  Configuration (c3,
%     -105..-135). Algorithms and Configuration share a bottom edge at -135;
%     the package boundary already sat at -138.5 because of Configuration, so
%     growing Algorithms costs NO page size. Its compartment follows the Tools
%     spacing, not Configuration's: box top -105, header -109, rule -113,
%     members from -114 (Configuration's rule is 9.5 mm down because it has a
%     two-line header, Algorithms' name is one line). Five member lines at
%     \scriptsize (2.82 mm each) end at ~-128.9, so the box keeps ~6 mm of
%     bottom padding, matching Configuration's 5.4 mm.
%     Measured member widths at \scriptsize, against 40 mm of text width:
%         $\sim$ assign_clusters        21.87 mm
%         $\sim$ representations        21.69 mm
%         $\sim$ segmentation           19.29 mm
%         \quad $\sim$ KMedoids, KMaxoids  32.45 mm  (widest)
%     Configuration's widest line is 39.63 mm and already fits, so nothing
%     here is close to the border.
%   externals y -143..-158  scikit-learn (c2)  Pyomo (c3)
%   legend rows y -171, -179, -187; row 2 has THREE columns (swatch at x 8,
%     dependency arrow 53..65, object flow arrow 98..110), which is what lets
%     the new shape swatch in without a fourth row. Measured label widths at
%     \scriptsize: "pipeline phase (action)" 28.08 mm ends at 43.6,
%     "\guillemotleft use\guillemotright dependency" 22.28 ends at 88.8, "object flow between phases"
%     33.13 ends at 144.6 — all inside the 157 right edge, so the drawing's
%     width is still set by the package boundary and does not move.
%   Lanes, each checked against every other line sharing it:
%     x = 150   Pipeline -> Results        (y -56 .. -42)
%     x = 104   Algorithms -> Pyomo        (y -120 .. -150.5); clear of
%               Configuration, which starts at x 109. It leaves the box at
%               -120, its vertical centre, NOT at the old -114: with a
%               compartment there, -114 is where the rule meets the border and
%               the edge would appear to spring out of the rule.
%     x =  77   API -> Pipeline (-48..-56), Pipeline -> Algorithms (-97..-105),
%               Algorithms -> scikit-learn (-135..-147); disjoint y-ranges.
%               The scikit-learn leg is 12 mm now that the box reaches -135,
%               down from 24, which is still ample for the arrowhead.
%   The "exact kmedoids" label hangs off the RIGHT side of the x = 104 lane,
%   anchored west at (105.5,-143.5), so it touches the edge it names and only
%   that one. Anchored east at (102,-142) it ran leftwards to x 82.6, which put
%   it 5.6 mm from the Algorithms -> scikit-learn lane at x = 77 and 1.5 mm
%   from its own: near enough to both to be read as labelling either. On the
%   right of the lane it is 28 mm clear of x = 77, 3.5 mm below the package
%   boundary at -138.5 and 2 mm above the Pyomo box at -147.
%
% ── TRAPS RE-CHECKED HERE (see the project notes) ─────────────────────────
%   * A TikZ node CLIPS NOTHING: an over-long line spills across its border and
%     silently sets the drawing's width. Every node below has a `text width`.
%   * `inner sep=0` on any full-width text block.
%   * Guillemets as \guillemotleft / \guillemotright — literal UTF-8 renders as
%     "nsharedz" under T1. \textperiodcentered{} for the middot, --- for the
%     em dash.
%   * A TikZ STYLE cannot be applied to part of a node's text; the greyed and
%     ambered stage lines use \newcommand DECLARATIONS instead.
%   * MEASURE text, do not estimate it. Widths that matter here, from
%     _measure.tex (\settowidth, then read out of the PDF — tectonic swallows
%     \typeout):
%         unstack_to_periods   \tiny roman   21.26 mm
%         unstack_to_periods   \tiny \texttt 16.80 mm
%         count occurrences    \tiny roman   18.76 mm
%         Representatives      \tiny roman   16.46 mm
%         exact kmedoids       \scriptsize   18.87 mm
%     A phase box is 22.5 mm wide with 19.5 mm of text width, so the roman form
%     overflowed the text box by 1.76 mm and cleared the border by ~0.6 mm,
%     which prints as touching it.
%   * The legend swatch takes `rounded corners=1.5mm`, not the phases' 2.2 mm.
%     A 2.2 mm radius on a 6 mm-tall swatch eats three quarters of its height
%     and reads as a lozenge rather than as the same shape, smaller.
%
% ── WHEN A STAGE NAME IS SET IN MONOSPACE ─────────────────────────────────
% The rule is: a stage written as a snake_case IDENTIFIER is set in \texttt,
% because it names a function; a stage written as prose is not, because it
% describes a step. `unstack_to_periods` and the activity's own name
% `run_pipeline` qualify. This is also what makes them fit — see the
% measurements above.
%
% Regenerate:
%   tectonic architecture_overview.tex
%   pdftocairo -svg architecture_overview.pdf architecture_overview.svg

\documentclass[10pt,border=1mm]{standalone}
\usepackage[T1]{fontenc}
\usepackage{lmodern}
\usepackage{tikz}
\usetikzlibrary{positioning,arrows.meta,fit,backgrounds,calc}

\definecolor{fzjblue}{HTML}{023D6B}
\definecolor{fzjdark}{HTML}{012845}
\definecolor{tintblue}{HTML}{DCE9F2}
\definecolor{phasefill}{HTML}{F4F8FB}
\definecolor{teal}{HTML}{21918C}
\definecolor{amber}{HTML}{A8681A}
\definecolor{extfill}{HTML}{E6E9EC}
\definecolor{extline}{HTML}{8A94A0}
\definecolor{greytext}{HTML}{3F4A55}

% DECLARATIONS, not TikZ styles: a style cannot be applied to part of a node's
% text, which is what the phase tags and the two kinds of stage line need.
\newcommand{\ph}{\tiny\color{greytext}}
\newcommand{\stg}{\tiny\color{fzjdark}}
\newcommand{\opt}{\tiny\color{amber}}
\newcommand{\kw}{\scriptsize\color{greytext}}

\begin{document}
\begin{tikzpicture}[
  x=1mm, y=1mm,
  api/.style={
    draw=fzjblue, line width=0.45mm, fill=tintblue, text=fzjdark,
    align=center, font=\small},
  int/.style={
    draw=fzjblue, line width=0.45mm, fill=white, text=fzjdark,
    align=center, font=\small},
  % UML Action: a ROUND-CORNERED rectangle. This is the only thing separating
  % the pipeline's temporal phases from the components around them.
  phase/.style={
    draw=fzjblue, line width=0.35mm, fill=phasefill, text=fzjdark,
    rounded corners=2.2mm,
    minimum width=22.5mm, minimum height=25mm, text width=19.5mm,
    align=center, font=\scriptsize},
  ext/.style={
    draw=extline, line width=0.4mm, fill=extfill, text=greytext,
    minimum width=40mm, minimum height=15mm, text width=36mm,
    align=center, font=\small},
  hdr/.style={font=\small, text=fzjdark, align=center},
  mem/.style={font=\scriptsize, text=fzjdark, align=left, anchor=north west},
  rule/.style={draw=fzjblue, line width=0.3mm},
  % UML Dependency: dashed line, open arrowhead. Keyword declared in the legend.
  dep/.style={-{Straight Barb[scale=1.6]}, dashed, draw=fzjblue, line width=0.3mm},
  extdep/.style={dep, draw=extline},
  % UML ObjectFlow between Actions: SOLID line, open arrowhead. The label names
  % the object that flows.
  conn/.style={-{Straight Barb[scale=2]}, draw=teal, line width=0.5mm},
  flbl/.style={font=\tiny, text=fzjdark, align=center, anchor=south,
               text width=14mm, inner sep=0.25mm},
  elbl/.style={font=\scriptsize, text=greytext, inner sep=1.4pt, fill=white},
  lgd/.style={font=\scriptsize, text=greytext, anchor=west}
]
\hyphenpenalty=10000 \exhyphenpenalty=10000

% ══ provided interface: the entry point, in UML ball notation ══════════════
\node[circle, draw=fzjblue, line width=0.4mm, fill=white, minimum size=4mm,
      inner sep=0] (iface) at (77,-6) {};
\draw[draw=fzjblue, line width=0.4mm] (iface.south) -- (77,-16);
\node[anchor=west, font=\scriptsize, text=fzjdark] at (80.5,-6) {aggregate()};

% ══ row A — the three components that expose public API ════════════════════
% The whole c1 column is left-aligned at x = 9, not x = 2, to open the lane the
% Tools -> Plotly dependency needs. See the LANE note in the header.
% unstack_to_periods sits here rather than in API. It IS public API
% (tsam.__all__), but its documented job is reshaping a series for a plotly
% heatmap (api.py:282) — it is never part of an aggregate() call. Grouping it
% with the other user-facing helpers says that; the cost is that the box is a
% purpose grouping here, since the function's module is api.py. Written without
% "(...)": at \scriptsize the parenthesised form sets 38.9 mm and would spill
% the 37 mm box, and `tuning` / `plot` carry no parens either.
\node[api, minimum width=41mm, minimum height=24mm] (tools) at (29.5,-28) {};
\node[hdr] at (29.5,-20) {Tools};
\draw[rule] (9,-24) -- (50,-24);
\node[mem, text width=37mm] at (11,-25)
  {+ tuning\\+ plot\\+ unstack\_to\_periods};

\node[api, minimum width=44mm, minimum height=28mm] (apic) at (77,-30) {};
\node[hdr] at (77,-20) {API};
\draw[rule] (55,-24) -- (99,-24);
\node[mem, text width=40mm] at (57,-25)
  {+ aggregate(data, n\_clusters,\\
   \quad cluster: ClusterConfig,\\
   \quad segments: SegmentConfig,\\
   \quad extremes: ExtremeConfig, \ldots)\\
   \quad $\rightarrow$ AggregationResult};

\node[api, minimum width=48mm, minimum height=26mm] (results) at (133,-29) {};
\node[hdr] at (133,-20) {Results};
\draw[rule] (109,-24) -- (157,-24);
\node[mem, text width=44mm] at (111,-25)
  {+ AggregationResult\\+ ClusteringResult\\+ AccuracyMetrics\\+ ConcurrencyMetrics};

% ══ Pipeline — a component whose second compartment holds the ACTIVITY ═════
\node[int, minimum width=148mm, minimum height=41mm] (pipeline) at (83,-76.5) {};
\node[hdr] at (83,-60) {Pipeline};
\draw[rule] (9,-64) -- (157,-64);

% The activity's name, in the upper left of the compartment that frames it —
% which is where UML puts it. It is the one place \guillemotleft activity\guillemotright appears; saying
% it here rather than in the legend is what tells the reader that the four
% round-cornered boxes below are steps of a run and not four more parts.
\node[anchor=north west, font=\tiny, text=greytext, inner sep=0.2mm]
  at (11,-64.5) {\guillemotleft activity\guillemotright{} \texttt{run\_pipeline}};

\node[phase] (p1) at ( 24.25,-80.5)
  {{\ph Phase 1}\\[0.4mm]\textbf{Prepare data}\\[0.8mm]
   {\stg normalize\\\texttt{unstack\_to\_periods}\\}{\opt apply weights\\add period sums}};
\node[phase] (p2) at ( 63.42,-80.5)
  {{\ph Phase 2}\\[0.4mm]\textbf{Cluster}\\[0.8mm]
   {\stg assign clusters\\representations}};
\node[phase] (p3) at (102.58,-80.5)
  {{\ph Phase 3}\\[0.4mm]\textbf{Refine}\\[0.8mm]
   {\opt extreme periods\\unweight\\}{\stg count occurrences\\}{\opt rescale\\segment}};
\node[phase] (p4) at (141.75,-80.5)
  {{\ph Phase 4}\\[0.4mm]\textbf{Build result}\\[0.8mm]
   {\stg denormalize\\reconstruct\\accuracy metrics}};

% Object flows between the actions, labelled with the object each conveys. The
% labels are free-standing nodes in the gaps, NOT pos= labels on the paths.
% A label node is 16 mm of text width plus 0.6 mm of inner sep a side = 17.2 mm,
% so the gap must exceed that: at the previous 24.45 mm parts the gap was
% 16.4 mm and "Refined Representatives" overlapped the Phase 4 box.
\draw[conn] ( 35.50,-80.5) -- ( 52.17,-80.5);
\draw[conn] ( 74.67,-80.5) -- ( 91.33,-80.5);
\draw[conn] (113.83,-80.5) -- (130.50,-80.5);
% Anchored south at -76.5, not -79: the flow sits at -80.5 and its Straight
% Barb at scale 2 reaches ~1.5 mm above the line, which was touching the
% descenders of the two-line labels. All three move together so the row stays
% level.
\node[flbl] at ( 43.83,-76.5)  {Prepared\\Data};
\node[flbl] at ( 83.00,-76.5) {Cluster\\Assignment};
% "Representatives" sets 17.4 mm at \tiny, wider than the 16.67 mm gap, so it
% overran the Phase 4 border. Hyphenation is off globally (\hyphenpenalty),
% so the break is made by hand.
\node[flbl] at (122.17,-76.5)  {Refined\\Represen-\\tatives};

% ══ row C ══════════════════════════════════════════════════════════════════
% Commons and Configuration carry \guillemotleft shared\guillemotright and NO arrows: they are used by
% every building block, and one arrow each would assert the wrong relationship.
% Commons has no API compartment, so the keyword rides in the node's own text
% and centres itself; a separate \node[hdr] sat visibly high in the box.
\node[int, minimum width=41mm, minimum height=18mm, text width=37mm]
  (commons) at (29.5,-114) {{\kw \guillemotleft shared\guillemotright}\\[0.3mm]Commons};

% Algorithms, with a members compartment. `~` is UML package visibility —
% visible inside tsam, not exported — which is the honest marker for a box the
% legend already colours "internal to tsam"; `+` would claim public API and `-`
% would deny that Pipeline calls these. The three named here are the three the
% phase boxes name as stages; the run-in line adds the two exact-optimisation
% classes, which is what the scikit-learn and Pyomo edges are about.
\node[int, minimum width=44mm, minimum height=30mm] (algorithms) at (77,-120) {};
\node[hdr] at (77,-109) {Algorithms};
\draw[rule] (55,-113) -- (99,-113);
\node[mem, text width=40mm] at (57,-114)
  {$\sim$ assign\_clusters\\$\sim$ representations\\$\sim$ segmentation\\
   \quad{\tiny exact methods}\\
   \quad $\sim$ KMedoids, KMaxoids};

% Distribution and MinMaxMean are in tsam.__all__ too, so the "+ = public API"
% convention was incomplete without them. They are not peers of the three
% configs — they are the values the `representation` attribute of ClusterConfig
% and SegmentConfig takes — hence the indented run-in line rather than two more
% "+" entries.
\node[api, minimum width=48mm, minimum height=30mm] (config) at (133,-120) {};
\node[hdr] at (133,-109.5) {{\kw \guillemotleft shared\guillemotright}\\[0.3mm]Configuration};
\draw[rule] (109,-114.5) -- (157,-114.5);
\node[mem, text width=44mm] at (111,-115.5)
  {+ ClusterConfig\\+ SegmentConfig\\+ ExtremeConfig\\
   \quad{\tiny .representation takes}\\
   \quad + Distribution, MinMaxMean};

% ══ the tsam package: UML package notation (body + folder tab) ═════════════
\begin{scope}[on background layer]
  \node[draw=fzjblue, line width=0.45mm, fill=white, inner sep=3.5mm,
        fit=(tools)(apic)(results)(pipeline)(commons)(algorithms)(config)(iface)]
        (tsambox) {};
  \draw[draw=fzjblue, line width=0.45mm, fill=white]
    (tsambox.north west) rectangle ++(30,8);
\end{scope}
\node[anchor=west, text=fzjblue, font=\small\bfseries]
  at ([shift={(2mm,4mm)}]tsambox.north west) {tsam};

% ══ externals ══════════════════════════════════════════════════════════════
% The pervasive dependencies go in the free c1 slot, with no edges, for the
% same reason Commons and Configuration have none: NumPy and pandas are used by
% every building block, so one arrow each would assert the wrong relationship.
% Plotly is not pervasive — only Tools uses it — and a full-width Pipeline
% whitebox blocks every vertical lane between row A and this row, so it cannot
% be given an honest edge either; it is named here with its user instead.
% Two boxes, not one: NumPy and pandas are \guillemotleft shared\guillemotright and Plotly is not, so
% putting them in a single box made the keyword claim something false about
% Plotly. The pocket below Tools (x 2..46, y -36..-56) is free but unusable —
% it is INSIDE the package rectangle, and an external dependency drawn there
% would read as part of tsam. A UML package is a rectangle; it cannot be
% notched around the box.
% Plotly is the UPPER box so the lane can reach its west edge without the
% vertical passing through the other box.
\node[ext, minimum width=41mm, minimum height=9mm, text width=37mm]
  (plotly) at ( 29.5,-150) {Plotly};
\node[ext, minimum width=41mm, minimum height=8mm, text width=37mm]
  (pydata) at ( 29.5,-160.5)
  {{\kw \guillemotleft shared\guillemotright}\\pandas \textperiodcentered{} NumPy};

\node[ext] (sklearn) at ( 77,-154.5) {scikit-learn};
\node[ext] (pyomo)   at (133,-154.5)
  {Pyomo\\{\tiny HiGHS \textperiodcentered{} CBC \textperiodcentered{} Gurobi
   \textperiodcentered{} CPLEX}};

% ══ dependencies — all orthogonal ══════════════════════════════════════════
\draw[dep] ( 50,-26)  -- ( 55,-26);                 % Tools -> API
\draw[dep] ( 99,-26)  -- (109,-26);                 % API -> Results
\draw[dep] ( 77,-44)  -- ( 77,-56);                 % API -> Pipeline
\draw[dep] (150,-56)  -- (150,-42);                 % Pipeline -> Results
\draw[dep] ( 77,-97)  -- ( 77,-105);                % Pipeline -> Algorithms

% Tools -> Plotly, down the only corridor that exists. Every column is
% left-aligned at x = 9 and the package boundary lands at 5.5, so the lane at
% x = 1 runs OUTSIDE the package with 4.5 mm of daylight — enough not to read
% as a second border. It leaves Tools' west edge, so it crosses the boundary
% once, at the top, which is exactly what a dependency leaving the package
% should do.
\draw[extdep] (  9,-26) -- (  1,-26) -- (  1,-150) -- (  9,-150);

\draw[extdep] ( 77,-135) -- ( 77,-147);             % Algorithms -> scikit-learn
\draw[extdep] ( 99,-120) -- (104,-120) -- (104,-154.5) -- (113,-154.5);
% On the RIGHT of its own lane, not the left. Anchored east at (102,-142) the
% label reached back to x 82.6, close enough to the scikit-learn lane at x = 77
% to be read as labelling that edge instead. Here it touches x = 104 with 1 mm
% of daylight and is 28 mm clear of x = 77.
\node[elbl, anchor=west] at (105.5,-143.5) {exact kmedoids};

% ══ key ════════════════════════════════════════════════════════════════════
% Row 1 colour swatches (what kind of component), row 2 shape and edge kinds
% (what kind of element), row 3 the two text conventions. Row 2 runs THREE
% columns so the round-cornered swatch fits without a fourth row, which the
% aspect budget cannot pay for.
\node[api, minimum width=12mm, minimum height=6mm] at ( 8,-171) {};
\node[lgd] at (15.5,-171) {exposes public API};

\node[int, minimum width=12mm, minimum height=6mm] at (52,-171) {};
\node[lgd] at (59.5,-171) {internal to tsam};

\node[ext, minimum width=12mm, minimum height=6mm, text width=8mm] at (95,-171) {};
\node[lgd] at (102.5,-171) {external dependency};

% The shape swatch. 1.5 mm of radius, not the phases' 2.2, because the swatch
% is 6 mm tall and 2.2 would round it into a lozenge.
\node[phase, rounded corners=1.5mm, minimum width=12mm, minimum height=6mm,
      text width=8mm] at (8,-179) {};
\node[lgd] at (15.5,-179) {pipeline phase (action)};

\draw[dep] (53,-179) -- (65,-179);
\node[lgd] at (66.5,-179)
  {\guillemotleft use\guillemotright{} dependency};

\draw[conn] (98,-179) -- (110,-179);
\node[lgd] at (111.5,-179) {object flow between phases};

\node[lgd, text=amber] at (4,-187) {\textit{amber}};
\node[lgd] at (17.5,-187) {stage runs only when configured};

\node[lgd, text=greytext] at (72,-187) {\guillemotleft shared\guillemotright};
\node[lgd] at (85.5,-187) {used by every building block};

\end{tikzpicture}
\end{document}
